% ---------
% --------
%!TEX TS-program = pdflatex
%!TEX encoding = UTF-8 Unicode
%!TEX spellcheck = English (Aspell)

\documentclass[11pt]{article}
\usepackage{lucy,handout,graphicx}
\usepackage{fourier-orns}
\usepackage{mathtools}
\usepackage[normalem]{ulem} % either use this (simple) or
\usepackage{tikz}

\def\Cdot{\mathbin{\text{\normalfont \textbullet}}}
\def\Sym#1{\texttt{\color{BrickRed}#1}}
\def\BlueSym#1{\textbf{\texttt{\color{RoyalBlue}#1}}}
\allowdisplaybreaks

\begin{document}

\begin{center}
\Large\textbf{CSCI 332, Fall 2026}%
\\
\LARGE\textbf{Homework 7 (Part 2---Part 1 is on Prairie Learn)}%
\\[0.5ex]
\large Due Monday, October 12 Anywhere on Earth (6am Tuesday)
\end{center}

\bigskip
\hrule
\bigskip

\subsection*{Submission Requirements}
\begin{itemize}
    \item Type or clearly hand-write your solutions into a PDF format so that they are legible and professional. Submit your PDF on Gradescope. 
    \item Do not submit your first draft. Type or clearly re-write your solutions for your final submission. If your submission is not legible, we will ask you to resubmit.
    \item Use Gradescope to assign problems to the correct page(s) in your solution. If you do not do this correctly, we will ask you to resubmit.
\end{itemize}

\subsection*{Academic Integrity}

Remember, you may access \EMPH{any} resource in preparing your solution to the homework. However, you \EMPH{must}
\begin{itemize}
    \item write your solutions in your own words, and
    \item credit every resource you use (for example: ``Bob Smith helped me on
    this problem. He took this course at UM in Fall 2020''; ``I found a solution
    to a problem similar to this one in the lecture notes for a different
    course, found at this link: www.profzeno.com/agreatclass/lecture10''; ``Here is a link to the full transcript of my chat with ChatGPT about the problems on this HW''.)
    Ask if you need any clarifications about how to properly cite your sources!
\end{itemize}




\newpage
%----------------------------------------------------------------------
    
\headers{CSCI 332}{Homework 7 (due October 12)}{Fall 2026}


\begin{enumerate}
%\parindent 1.5em \itemsep 4ex plus 0.5fil

%----------------------------------------------------------------------

\item (0 points) List the resources you used for each problem in this homework. If you used no resources except
lecture notes and videos and the linked textbook, you can say ``none''. If you used
an AI tool, you should put a link to your full conversation here, and you must
use the
\href{https://umaal.org/files/teaching/csci332_fall2026/agents.md}{provided
prompt}
 to guide the AI to act
as a tutor, not an answer-provider.

You will be asked to resubmit your homework if this section is blank.

Remember that no matter what resources you consulted, you must write your solutions
\emph{in your own words}.

\item (5 points)

This written homework will extend on the PrairieLearn assignment's final problem:

You are given a sequence of digits $A[1..n]$ where each digit is between 1 and 9
(inclusive). You are asked to insert $+$ signs in between the digits to partition them
into terms which will be added together. The length of each term in the summation
must be either 1 or 2 digits.

What is the maximum sum that can be achieved under these constraints?

For example, if $A=[1,9,1,2]$, the maximum sum is 1 + 91 + 2 = 94.

\begin{enumerate}

\item  The PrairieLearn assignment asked you to write a viable
subproblem definition that can be used to solve this problem. Copy your
subproblem definition here.

\item  Write a recursive definition for your subproblem. Make sure to
include base case(s) as needed!

\item  Describe how you will memoize your recursive solution. What
data structure will you use and in what order will you fill it?

\item Write pseudocode for a dynamic programming algorithm that
uses your subproblem definition to solve the problem.

\item  Implement your pseudocode in the programming language of your
choice and use your implementation to verify that your algorithm is correct.
This means that you should run your code on a few examples.  Copy your code and
the output on your examples here.

\end{enumerate}
\end{enumerate}
%%%%%%%%%%%%%%%%%%%%
\end{document}
